\clearpage
\VerdierPhysicalPageAnchor{0138}%
\begingroup
\small
\setlength{\emergencystretch}{2em}%
\setlength{\parskip}{0pt}%
\setlength{\abovedisplayskip}{1pt plus 1pt minus 1pt}%
\setlength{\belowdisplayskip}{1pt plus 1pt minus 1pt}%

\begingroup\itshape
\VerdierSemanticAnchor{ch2-2-2-4-b}
\noindent b) 当范畴 $\mathcal C$ 是加性范畴时，范畴
$\mathcal C(S^{-1})$ 是加性范畴，且函子 $Q$ 是加性函子。

\VerdierSemanticAnchor{ch2-2-2-4-c}
\par\noindent c) 设 $\widehat S$ 为 $\mathcal C$ 中被 $Q$ 送到
$\mathcal C(S^{-1})$ 的同构的态射集合。
\VerdierSemanticAnchor{ch2-2-2-4-c-saturation}%
集合 $\widehat S$ 是 $\mathcal C$ 的饱和乘法系。$\mathcal C$ 中每个
包含 $S$ 的饱和乘法系都包含 $\widehat S$。
\endgroup

\VerdierSemanticAnchor{ch2-2-2-4-proof-b}
\par\noindent 第一个断言立即由 $\mathcal C(S^{-1})$ 的泛性质得到：
范畴 $\bigl(\mathcal C^\circ(S^{-1})\bigr)^\circ$ 是该泛问题的一个解。
第二个断言立即由定理 (2.2.2) 的断言 b) 与 c) 得到。
\VerdierSemanticAnchor{ch2-2-2-4-proof-c-characterization}%
由 $\mathcal C(S^{-1})$ 中态射的描述立即可知，$\mathcal C$ 的态射
$f$ 被 $Q$ 送到同构，当且仅当 $\mathcal C$ 中存在两个态射 $g$ 与
$h$，使得 $gf$ 与 $fh$ 有定义并且属于 $S$。
\VerdierSemanticAnchor{ch2-2-2-4-proof-c-sm2-sm3}%
由此立即可知，集合 $\widehat S$ 具有性质 SM2) 与 SM3)。
\VerdierSemanticAnchor{ch2-2-2-4-proof-c-sm1-sm4}%
此外，由 $\widehat S$ 的定义可知，该集合具有性质 SM1) 与 SM4)。
\VerdierSemanticAnchor{ch2-2-2-4-proof-c-conclusion}%
因此，$\widehat S$ 是 $\mathcal C$ 的饱和乘法系。(2.2.4) 的最后一个
断言是显然的。

\VerdierSemanticAnchor{ch2-2-2-5}
\smallskip
\noindent\textbf{2.2.5.} 设 $\mathcal C$ 为范畴，$S$ 为
$\mathcal C$ 的乘法系。
\VerdierSemanticAnchor{ch2-2-2-5-localized-category}%
范畴 $\mathcal C(S^{-1})$ 称为 $\mathcal C$ 关于乘法系 $S$ 的
\emph{局部化范畴}。
\VerdierSemanticAnchor{ch2-2-2-5-localization-functor}%
函子 $Q\colon\mathcal C\longrightarrow\mathcal C(S^{-1})$ 称为
\emph{局部化函子}。
\VerdierSemanticAnchor{ch2-2-2-5-saturation}%
乘法系 $\widehat S$ 称为乘法系 $S$ 的\emph{饱和化}。

\VerdierSemanticAnchor{ch2-2-2-6}
\smallskip
\noindent\textbf{定理 2.2.6.}\quad
\begingroup\itshape
\VerdierSemanticAnchor{ch2-2-2-6-setup}%
设 $\mathcal D$ 为三角范畴，$S$ 为 $\mathcal D$ 的乘法系 (2.1.4)。
以 $\mathcal D(S^{-1})$ 表示 $\mathcal D$ 的局部化范畴，以
$Q\colon\mathcal D\longrightarrow\mathcal D(S^{-1})$ 表示局部化函子。

\VerdierSemanticAnchor{ch2-2-2-6-a}
\par\noindent a) 在 $\mathcal D(S^{-1})$ 上存在唯一的严格
$\mathbb Z$-范畴结构，使 $Q$ 为严格 $\mathbb Z$-函子。

\VerdierSemanticAnchor{ch2-2-2-6-b}
\par\noindent b) 在 $\mathcal D(S^{-1})$ 上存在唯一的三角结构，使
$Q$ 为正合函子。该结构的特异三角，是那些同构于 $\mathcal D$ 的
特异三角在 $Q$ 下的像的三角。

\VerdierSemanticAnchor{ch2-2-2-6-c-factorization}
\par\noindent c) 到某个三角范畴（分别地，阿贝尔范畴）、并把 $S$
中态射送到同构的每个正合（分别地，上同调）函子
$F\colon\mathcal D\longrightarrow\mathcal D'$，都唯一地通过
$Q\colon\mathcal D\longrightarrow\mathcal D(S^{-1})$ 分解为正合
（分别地，上同调）函子
$G\colon\mathcal D(S^{-1})\longrightarrow\mathcal D'$。
\VerdierSemanticAnchor{ch2-2-2-6-c-morphism-lifting}%
设 $F_1,F_2\colon\mathcal D\rightrightarrows\mathcal D'$ 为两个把
$S$ 中态射送到同构的正合（分别地，上同调）函子；设
$G_1,G_2\colon\mathcal D(S^{-1})\rightrightarrows\mathcal D'$ 为它们
各自的分解，并设 $m\colon F_1\longrightarrow F_2$ 为正合（分别地，
上同调）函子的态射。存在唯一的函子态射
$p\colon G_1\longrightarrow G_2$，使 $m=p\star Q$。
\endgroup

\VerdierSemanticAnchor{ch2-2-2-6-proof-a}
\par\noindent 第一个断言立即由定理 (2.2.2) 与 SM5) 得到。
\VerdierSemanticAnchor{ch2-2-2-6-proof-c-p139-boundary}%
第三个断言立即由

\VerdierSemanticAnchor{ch2-2-2-6-p139-boundary}
\endgroup
